\section{Contributing to xeus-haskell}
\label{sec:contributing}


Thank you for your interest in contributing to \texttt{xeus-haskell}!
Xeus and xeus-haskell are subprojects of Project Jupyter and subject to the\href{https://github.com/jupyter/governance}{Jupyter governance} and\href{https://github.com/jupyter/governance/blob/master/conduct/code\_of\_conduct.md}{Code of conduct}.
\subsection{General Guidelines}


For general documentation about contributing to Jupyter projects, see the\href{https://jupyter.readthedocs.io/en/latest/contributor/content-contributor.html}{Project Jupyter Contributor Documentation}.
\subsection{Community}


The Xeus team organizes public video meetings. The schedule for future meetings andminutes of past meetings can be found on our\href{https://jupyter-xeus.github.io/}{team compass}.
\subsection{Getting Started}


\subsubsection{Prerequisites}


The project uses \href{https://pixi.sh/}{Pixi} for dependency management and workflow automation. If you haven't installed it yet:
\begin{minted}{text}
curl -fsSL https://pixi.sh/install.sh | sh
\end{minted}


\subsubsection{Setting Up Your Environment}


\begin{enumerate}
\item Fork the repository on GitHub.
\item Clone your fork locally:
\end{enumerate}


\begin{minted}{text}
git clone https://github.com/<your-username>/xeus-haskell.git
cd xeus-haskell
\end{minted}


\begin{enumerate}
\item Initialize the development environment:
\end{enumerate}


\begin{minted}{text}
pixi run -e default prebuild
\end{minted}


The \texttt{default} environment includes all necessary tools (CMake, C++ compilers, Python for testing, etc.).
\subsection{Development Workflow}


\subsubsection{Building the Kernel}


To build and install the kernel in your local environment:
\begin{minted}{text}
pixi run -e default build
pixi run -e default install
\end{minted}


\subsubsection{Running Tests}


We use \texttt{pytest} and \texttt{ctest} for verifying the kernel:
\begin{minted}{text}
# Run C++ tests
pixi run -e default ctest

# Run Python-based Jupyter kernel tests
pixi run -e default pytest
\end{minted}


\subsubsection{Working with WebAssembly (JupyterLite)}


If you are contributing to the WebAssembly build, you need the \texttt{wasm-build} and \texttt{wasm-host} environments:
\begin{minted}{text}
# Prepare the WASM host environment
pixi install -e wasm-build
pixi install -e wasm-host

# Build for WebAssembly
pixi run -e wasm-build prebuild
pixi run -e wasm-build build
pixi run -e wasm-build install

# Run the JupyterLite server locally
pixi run -e wasm-build serve
\end{minted}


\subsection{Project Structure}


\begin{itemize}
\item \texttt{src/}: C++ source code for the kernel and REPL bridge.
\item \texttt{include/}: C++ header files.
\item \texttt{share/}: Haskell side of the implementation, including \texttt{Repl.lhs} and MicroHs libraries.
\item \texttt{test/}: Python tests for kernel functionality.
\item \texttt{docs/}: LaTeX source files for project documentation.
\item \texttt{xhaskell.tex}: Root LaTeX document that assembles docs and literate Haskell sources.
\end{itemize}


\subsection{Style Guidelines}


\begin{itemize}
\item \textbf{C++}: Follow modern C++ practices. We use Xeus as the base library.
\item \textbf{Haskell}: Since we use MicroHs, avoid dependencies or language features not supported by the MicroHs compiler.
\item \textbf{Documentation}: Write documentation in LaTeX (\texttt{.tex}) and keep examples runnable with the current Pixi workflows.
\end{itemize}

\subsubsection{Building Documentation}

Generate the project PDF documentation with:
\begin{minted}{text}
pixi run -e docs generate
\end{minted}


\subsection{Submitting Changes}


\begin{enumerate}
\item Create a new branch for your feature or bugfix.
\item Ensure all tests pass.
\item Submit a Pull Request with a clear description of the changes.
\end{enumerate}


\subsection{Community}


Join our public video meetings! The schedule and minutes are available on our \href{https://jupyter-xeus.github.io/}{Team Compass}.
